
Aligning language models to both aggregate human preferences (AI-to-Human alignment) and user-specific controls (Human-to-AI alignment) typically requires sequential training stages that risk catastrophic forgetting when the second objective degrades the first. We demonstrate that Direct Preference Optimization (DPO) and attribute conditioning can be jointly optimized in a single training run through gradient-compatible multi-task learning. Gradient angle measurements between DPO and attribute loss gradients averaged 78.5° (SD: 12.8°) across 100 proof-of-concept training steps, with zero measurements exceeding the 120° catastrophic interference threshold established in multi-task learning theory. Both DPO loss and attribute loss decreased monotonically (5.8\% and 21.3\% reduction respectively), confirming convergence without destructive task conflict.

The jointly trained model achieved 54.07\% preference win rate (approximately 94\% retention of the 57.5\% standalone DPO baseline) and 65.14\% attribute steering accuracy (45 percentage points above the 20\% random chance baseline on 5-level classification), both exceeding proof-of-concept feasibility thresholds. Linear probing analysis revealed that the joint model encodes preference information with 100\% classification accuracy from hidden states, demonstrating that multi-task training creates shared representations satisfying both objectives. However, attribute encoding analysis failed due to synthetic label contamination ($R^{2}=-1.324)$, and comparison against sequential training baselines was not conducted.

These findings establish gradient compatibility as a quantitative design principle for multi-objective language model alignment. The observed <120° gradient angle criterion provides practitioners with a falsifiable predictor of joint training feasibility, applicable beyond the specific DPO-attribute combination tested here to other multi-objective alignment scenarios including Constitutional AI constraints, safety objectives, and capability preservation.
